% !TEX root = ../thesis.tex

% Abstract -- German (Zusammenfassung) + English

\begin{otherlanguage}{ngerman}
    \section*{Zusammenfassung}

    Die Freigabe medizinischer Daten f\"ur die Forschung wird durch
    Datenschutzregelungen wie HIPAA und DSGVO eingeschr\"ankt, w\"ahrend sich
    klassische Anonymisierungsverfahren gegen\"uber Verkn\"upfungs- und
    Re-Identifikationsangriffen wiederholt als unzureichend erwiesen haben.
    Die Erzeugung synthetischer Daten bietet eine Alternative, doch der
    Zielkonflikt zwischen Nutzbarkeit und Privatsph\"are ist \"uber
    grundlegend verschiedene generative Paradigmen hinweg bislang nur
    unzureichend charakterisiert. Diese Arbeit pr\"asentiert einen
    einheitlichen empirischen Vergleich von f\"unf Generatoren auf dem
    Datensatz \emph{Diabetes 130-US Hospitals}, der die GAN- und die
    LLM-Familie jeweils mit und ohne formale Differential Privacy abdeckt:
    HealthGAN, PATE-GAN, Pythia-70M-LoRA, Pythia-1B-LoRA und
    Pythia-1B-DP-LoRA. Jeder Generator wird anhand skalarer Metriken entlang
    zweier S\"aulen bewertet: \emph{Utility}, die statistische \"Ahnlichkeit
    zu den realen Daten mit der pr\"adiktiven Leistung nachgelagerter, auf
    synthetischen Datens\"atzen trainierter Modelle verbindet, und
    \emph{Privacy}, die instanzbasierte Distanzma{\ss}e sowie Membership- und
    Attribute-Inference-Risiken umfasst.

    Keine der beiden Familien dominiert durchg\"angig. Ohne formale Garantie
    bleibt die Wahl ein echter Zielkonflikt: HealthGAN erreicht die beste
    marginale Fidelit\"at und als einziger nutzbarer Generator ein
    N\"achste-Nachbarn-Distanzverh\"altnis an der Einheitsreferenz. Die
    Sprachmodelle erzielen die h\"ochste pr\"adiktive Utility, wobei
    Pythia-1B der realen Referenz am n\"achsten kommt (AUC $0.682$
    gegen\"uber $0.701$). Unter einem angeglichenen Budget
    ($\varepsilon = 5.0$) verschiebt sich die Balance zugunsten der
    LLM-Familie: Pythia-1B-DP bleibt f\"ur einen R\"uckgang von etwa $0.10$
    im TSTR-AUC-ROC nutzbar, beseitigt zugleich das beim nicht-privaten
    Gegenst\"uck beobachtete Attribute-Inference-Leck und senkt die
    Outlier-Linkability auf null. Das hier evaluierte differentiell private
    GAN bricht dagegen auf beiden Utility-Zweigen ein und erzielt lediglich
    die scheinbare Privatheit einer von der realen Datenverteilung
    entfernten Ausgabe. Die Kosten der Differential Privacy sind somit stark
    familienabh\"angig. Die Arbeit liefert eine reproduzierbare
    Evaluationspipeline, einen direkten Benchmark entlang der GAN-LLM- und
    der DP-nicht-DP-Achse sowie quantitative Evidenz zu den Grenzkosten
    formaler Differential Privacy innerhalb jeder Generatorfamilie.
\end{otherlanguage}
\cleardoublepage

\begin{otherlanguage}{english}
    \section*{Abstract}

    Privacy regulations such as Health Insurance Portability and Accountability Act 
    (HIPAA) and General Data Protection Regulation (GDPR) provide the governance 
    framework within which medical data can be shared and reused responsibly 
    for research., while classical anonymisation has
    proven insufficient under linkage and re-identification attacks. Synthetic
    data generation offers an alternative, yet the trade-off between utility and
    privacy remains poorly characterised across fundamentally different
    generative paradigms. This thesis presents a unified empirical comparison of
    five generators on the \emph{Diabetes 130-US Hospitals} dataset, spanning
    the Generative Adversarial Network (GAN) and Large Language Model (LLM) 
    families both with and without formal differential privacy:
    HealthGAN, PATE-GAN, Pythia-70M-LoRA, Pythia-1B-LoRA, and Pythia-1B-DP-LoRA.
    Every generator is assessed with single-number metrics across two pillars:
    \emph{utility}, combining statistical similarity to the real data with the
    downstream predictive performance of models trained on the synthetic
    records, and \emph{privacy}, covering instance-level distance and
    membership- and attribute-inference risk.

    Without a formal guarantee the choice is a
    genuine trade-off: HealthGAN attains the best marginal fidelity and the only
    distance ratio at the nearest-neighbour unity reference. The language models
    reach the highest predictive utility, with Pythia-1B closest to the
    real-data ceiling (AUC $0.682$ against $0.701$). Under a matched
    $(\varepsilon = 5.0)$ budget the balance tips toward the LLM family:
    Pythia-1B-DP stays usable for roughly a $0.10$ drop in TSTR AUC-ROC while
    removing its non-private counterpart's observed attribute-inference leak
    and driving outlier linkability to zero. The differentially private GAN
    evaluated here, by contrast, collapses on both utility branches and earns
    only the apparent privacy of off-manifold output. The cost of differential privacy is thus
    strongly family-dependent. The thesis contributes a reproducible evaluation
    pipeline, a head-to-head benchmark spanning the GAN-LLM and DP-non-DP axes,
    and quantitative evidence on the marginal cost of formal differential
    privacy within each generative family.
\end{otherlanguage}
\cleardoublepage
